\chapter*{Sažetak}
\addcontentsline{toc}{chapter}{Sažetak}

Obveznik u Hrvatskoj od 2026.\ mora svaki eRačun dostaviti primatelju i prijaviti
ga Poreznoj upravi, i to dvama kanalima koje propis ne povezuje. Rad izvodi načine
na koje se taj tijek može pokvariti. Izvod se oslanja na Zakon, Pravilnik i
tehničku specifikaciju, pa je svaka tvrdnja provjerljiva bez pristupa ijednom
sustavu.

Izvedeno je jedanaest načina otkazivanja. Jedan od njih pokazuje da način
identifikacije ispravljenog eRačuna nije potpuno određen. Ispravak bez poreznog
učinka izdaje se pod istim brojem, provjera S008 duplikat prepoznaje po složenom
identifikatoru koji ne sadrži indikator kopije, a izvori ne određuju mora li
ispravak zadržati datum izvornika. Zadržavanje datuma po opisanom ključu vodi u
odbijanje, a novi datum uz referencu na prethodni račun tu provjeru izbjegava.
Rad zato utvrđuje nedorečenost sučelja; ne tvrdi da je svaki ispravak nužno
odbijen ni da bi stvarni sustav drugu granu prihvatio.

Taj slučaj nije jedini. Dio posljedica nastaje iz svojstava raspodijeljene
razmjene, a dio iz obvezne arhitekture ili iz pravila koja ne određuju semantiku
kvara. Rad za drugu skupinu upotrebljava radni pojam \emph{regulatorno inducirano
ograničenje}. Pokrivenost je ocijenjena matricom jedanaest načina otkazivanja i
šest obrazaca, uz označenu razinu dokaza za svaku čeliju. Obrasci ne uklanjaju
nijedan način otkazivanja, ali sužavaju prozor nepoznatog ishoda ili taj ishod
čine vidljivim. Njihov učinak ovisi o granici i konfiguraciji: isti mehanizam
može zaštititi lokalne resurse, ali loše podešen može smanjiti vrijeme preostalo
do isteka zakonskog roka.

\vspace{1em}
\noindent\textbf{Ključne riječi:} tolerancija na greške, raspodijeljeni sustavi,
fiskalizacija eRačuna, obrasci otpornosti, regulatorno inducirano ograničenje,
mjerilo izvedeno iz propisa

\chapter*{Abstract}
\addcontentsline{toc}{chapter}{Abstract}

From 2026, a Croatian business subject to fiscalization must deliver every
eInvoice to its recipient and report it to the Tax Administration, through two
channels that the regulation does not connect. This thesis derives the ways in
which that flow can break. The derivation draws on the Act, the Ordinance and
the technical specification, so every claim is verifiable without access to any
system.

Eleven failure modes are derived. One of them shows that the identification of a
corrected eInvoice is not fully specified. A correction without tax effect uses
the same invoice number, while the S008 duplicate check uses a composite
identifier that excludes the copy indicator, and the sources do not determine
whether the correction retains the original issue date. Retaining the date leads to
rejection under the documented key, while a new date with a reference to the
preceding invoice avoids that check. The thesis therefore identifies an
interface ambiguity; it neither claims that every correction is necessarily
rejected, nor that the production system would accept the other branch.

That case is not the only one. Some consequences arise from distributed
exchange, while others arise from the mandated architecture or from rules that
leave failure semantics undefined. The thesis uses
\emph{regulation-induced constraint} as a working term for the latter group. Coverage is assessed with a matrix of eleven failure modes and six
patterns, with a labelled evidence level for every cell. No pattern removes a
failure mode, but patterns narrow the window of an unknown outcome or make that
outcome observable. Their effect depends on the boundary and configuration: the
same mechanism may protect local resources, while a poor configuration may
reduce the time remaining before a statutory deadline.

\vspace{1em}
\noindent\textbf{Keywords:} fault tolerance, distributed systems, eInvoice
fiscalization, resilience patterns, regulation-induced constraint,
regulation-derived oracle
